\subsection{Flag quotient, tangent decomposition, and retractions}
The geometric variable is a chain of nested subspaces. Rotating a basis inside
one block leaves that chain unchanged, which motivates the quotient and its
horizontal tangent space. The orthogonal block decomposition also separates
captured gradient energy from the residual exactly. Coordinatewise adaptive
moments require additional information about the chosen basis; their behavior
is examined in Proposition~\ref{prop:tf-gauge}.

\begin{definition}[Ambient spaces and block selectors]
\begin{gather}
 \langle X,Y\rangle_F:=\tr(X^\top Y),\quad
 \|X\|_F^2:=\langle X,X\rangle_F,\quad
 \|X\|_{\op}:=\sup_{\|v\|_2=1}\|Xv\|_2,\\
 \sym A:=\tfrac12(A+A^\top),\qquad
 \skw A:=\tfrac12(A-A^\top),\\
 0=r_0<r_1<\cdots<r_K=R<d\le m,\quad
 q_j:=r_j-r_{j-1},\quad q_{K+1}:=d-R,\\
 I_{K+1}:=\{R+1,\ldots,d\},\\
 I_j:=\{r_{j-1}+1,\ldots,r_j\},\quad
 E_j:=\diag(\mathbf1_{I_j})\in\RR^{R\times R},\quad
 F_k:=\sum_{j\le k}E_j,\\
 \St(d,R):=\{U\in\RR^{d\times R}:U^\top U=I_R\},\\
 \mathcal G:=\prod_{j=1}^K\mathrm O(q_j),\qquad
 U\sim V\ \Longleftrightarrow\ 
 V=U\diag(O_1,\ldots,O_K),\quad O_j\in\mathrm O(q_j),\\
 \mathfrak F(q_1,\ldots,q_{K+1})
 :=\St(d,R)/\mathcal G
 \simeq\mathrm O(d)/\prod_{j=1}^{K+1}\mathrm O(q_j),\\
 U_j:=U_{:,I_j},\quad
 P_k(U):=UF_kU^\top,\quad
 \Pi_j(U):=UE_jU^\top,\quad \Pi_{K+1}(U):=I_d-UU^\top.
 \label{eq:tf-projectors}
\end{gather}
\end{definition}
\begin{proposition}[Projector algebra]\label{lem:tf-energy}
\begin{gather}
 P_k^\top=P_k=P_k^2,\quad \tr P_k=r_k,\quad
 P_kP_l=P_{\min(k,l)},\\
 \Pi_i^\top=\Pi_i,\quad \Pi_i\Pi_j=\delta_{ij}\Pi_i,
 \quad\sum_{j=1}^{K+1}\Pi_j=I_d,\\
 G=\sum_{j=1}^{K+1}G\Pi_j,\quad
 \langle G\Pi_i,G\Pi_j\rangle_F=0\quad(i\ne j),\\
 \|G\|_F^2=\sum_{j=1}^K\|GU_j\|_F^2+
 \|G\Pi_{K+1}\|_F^2,\\
 \|G-GP_k\|_F^2=\tr(G^\top G)-\tr(P_kG^\top G).
\end{gather}
\end{proposition}
\begin{proof}
\begin{align}
 P_kP_l&=UF_k(U^\top U)F_lU^\top=UF_{\min(k,l)}U^\top,\\
 \Pi_i\Pi_j&=UE_iE_jU^\top=\delta_{ij}UE_iU^\top,\quad i,j\le K,\\
 (I-UU^\top)UE_j&=0,\qquad
 \sum_{j\le K}UE_jU^\top=UU^\top,\\
 \langle G\Pi_i,G\Pi_j\rangle_F
 &=\tr(\Pi_iG^\top G\Pi_j)
 =\tr(\Pi_j\Pi_iG^\top G)=0,\\
 \|GU_jU_j^\top\|_F^2
 &=\tr(U_jU_j^\top G^\top GU_jU_j^\top)
 =\tr(U_j^\top G^\top GU_j),\\
 \|G(I-P_k)\|_F^2
 &=\tr((I-P_k)G^\top G(I-P_k))
 =\tr(G^\top G)-\tr(P_kG^\top G).
\end{align}
\end{proof}
\begin{proposition}[Vertical and horizontal spaces]\label{prop:rigid-horizontal}
\begin{align}
 T_U\St(d,R)&=\{Z:U^\top Z+Z^\top U=0\},\\
 \Ver_U&=\{U\Omega:\Omega^\top=-\Omega,
                  \ \Omega=\blkdiag(\Omega_{11},\ldots,\Omega_{KK})\},\\
 \Hor_U&=\{Z\in T_U\St(d,R):E_jU^\top ZE_j=0\ (1\le j\le K)\},\\
 T_U\St(d,R)&=\Hor_U\mathbin{\perp}\Ver_U,\\
 \mathsf H_U(A)&=A-U\sym(U^\top A)
   -U\blkdiag\bigl(\skw(U^\top A)\bigr),
 \label{eq:rigid-horizontal}\
 \dim\mathfrak F&=dR-\frac{R^2+\sum_{j=1}^Kq_j^2}{2}.
\end{align}
\end{proposition}
\begin{proof}
\begin{align}
 0&=\left.\frac{\mathrm d}{\mathrm ds}(U+sZ)^\top(U+sZ)\right|_{s=0}
   =U^\top Z+Z^\top U,\\
 Z&=U\Omega+U_\perp B,\quad
 \Omega=U^\top Z=-\Omega^\top,\quad B=U_\perp^\top Z,\\
 \langle Z,U\Lambda\rangle_F
 &=\tr(\Omega^\top\Lambda)
 =\sum_j\tr(\Omega_{jj}^\top\Lambda_{jj}),\\
 Z\perp\Ver_U&\Longleftrightarrow\Omega_{jj}=0\quad(\forall j),\\
 A&=\underbrace{U\sym(U^\top A)}_{(T_U\St)^\perp}
 +\underbrace{U\blkdiag(\skw(U^\top A))}_{\Ver_U}
 +\underbrace{\mathsf H_U(A)}_{\Hor_U},\\
 \langle U S,U\Omega\rangle_F&=\tr(S\Omega)=0
 \quad(S^\top=S,\ \Omega^\top=-\Omega),\\
 \dim\Hor_U&=dR-\frac{R(R+1)}2-\sum_j\frac{q_j(q_j-1)}2.
\end{align}
\end{proof}
We use projector distances to compare subspaces independently of their basis
representatives. The aligned-basis distance is convenient for moment transport;
the following comparisons allow estimates in these two representations to be
combined with explicit constants.

\begin{definition}[Quotient metric and projector distances]
\begin{gather}
 g_{[U]}(\xi,\eta):=\langle\xi_U^{\Hor},\eta_U^{\Hor}\rangle_F,\\
 d_P(U,V)^2:=\sum_{k=1}^K\|P_k(U)-P_k(V)\|_F^2,
 \label{eq:tf-distance}\\
 d_B(U,V)^2:=\sum_{j=1}^K\min_{O_j\in\mathrm O(q_j)}
                    \|U_j-V_jO_j\|_F^2,\\
 d_\Pi(U,V)^2:=\sum_{j=1}^K\|\Pi_j(U)-\Pi_j(V)\|_F^2,\\
 \overline d_\Pi(U,V)^2:=\sum_{j=1}^{K+1}\|\Pi_j(U)-\Pi_j(V)\|_F^2.
\end{gather}
\end{definition}
\begin{lemma}[Overlap transport identity and sharp metric comparisons]\label{lem:rigid-metrics}
For $T_{ja}:=\tr(\Pi_j(U)\Pi_a(V))$, $1\le j,a\le K+1$, set
$W:=\sum_{j\ne a}T_{ja}$. Then
\begin{gather}
 d_P^2=\sum_{j,a}|j-a|T_{ja},\quad\overline d_\Pi^2=2W,\\
 d_B^2\le d_\Pi^2\le2d_B^2,\qquad
 d_\Pi^2\le\overline d_\Pi^2\le2d_P^2,\\
 d_P^2\le\frac K2\overline d_\Pi^2\le Kd_\Pi^2\le2K d_B^2.
\end{gather}
The constants in $\overline d_\Pi^2\le2d_P^2$ and
$d_P^2\le(K/2)\overline d_\Pi^2$ are attained by adjacent-block and
first-to-residual rotations, respectively.
\end{lemma}
\begin{proof}
\begin{align}
 U_j^\top V_j&=L_j\diag(\sigma_{j1},\ldots,\sigma_{jq_j})R_j^\top,
 \quad0\le\sigma_{ji}\le1,\\
 \min_{O_j}\|U_j-V_jO_j\|_F^2&=2\sum_i(1-\sigma_{ji}),\\
 \|\Pi_j(U)-\Pi_j(V)\|_F^2&=2\sum_i(1-\sigma_{ji}^2),\\
 1-\sigma&\le1-\sigma^2\le2(1-\sigma),\quad\sigma\in[0,1],\\
 \sum_aT_{ja}&=q_j,\quad\sum_jT_{ja}=q_a,\\
 \|P_k(U)-P_k(V)\|_F^2
 &=\sum_{j\le k<a}T_{ja}+\sum_{a\le k<j}T_{ja},\\
 \sum_k\|P_k(U)-P_k(V)\|_F^2&=\sum_{j,a}|j-a|T_{ja},\\
 \sum_j\|\Pi_j(U)-\Pi_j(V)\|_F^2&=2\sum_{j\ne a}T_{ja}=2W,\\
 \|\Pi_{K+1}(U)-\Pi_{K+1}(V)\|_F^2
 &=2\sum_{a\le K}T_{K+1,a}\le W.
\end{align}
Thus $W\le d_P^2\le KW$ and $d_\Pi^2\ge W$. A two-plane rotation between
blocks $a<b$ has $W=2\sin^2\theta$ and $d_P^2=(b-a)W$, proving the equality
claims. Keeping the overlap flow avoids repeated triangle inequalities for
the nested prefixes.
\end{proof}
\begin{lemma}[Infinitesimal projector metric]\label{lem:rigid-differential}
For $Q=[U,U_\perp]\in\mathrm O(d)$, $\Omega^\top=-\Omega$,
$\Omega_{jj}=0$, and $U(s)=(Qe^{s\Omega})_{:R}$,
\begin{align}
 \dot P_k(0)&=Q[\Omega,\widehat F_k]Q^\top,\quad
 \widehat F_k:=\diag(I_{r_k},0_{d-r_k}),\\
 \sum_k\|\dot P_k(0)\|_F^2
 &=2\sum_{1\le a<b\le K+1}(b-a)\|\Omega_{ab}\|_F^2,\\
 \|\dot U(0)\|_F^2
 &=2\sum_{a<b\le K}\|\Omega_{ab}\|_F^2
   +\sum_{a\le K}\|\Omega_{a,K+1}\|_F^2,\\
 \|\dot U(0)\|_F^2
 &\le\sum_k\|\dot P_k(0)\|_F^2
 \le2K\|\dot U(0)\|_F^2.
\end{align}
\end{lemma}
\begin{proof}
\begin{align}
 P_k(s)&=Qe^{s\Omega}\widehat F_ke^{-s\Omega}Q^\top,\\
 [\Omega,\widehat F_k]_{ab}
 &=(\mathbf1_{b\le k}-\mathbf1_{a\le k})\Omega_{ab},\\
 \|[\Omega,\widehat F_k]\|_F^2
 &=2\sum_{a\le k<b}\|\Omega_{ab}\|_F^2,\\
 \sum_{k=1}^K\mathbf1_{a\le k<b}&=b-a,\\
 \|(Q\Omega)_{:R}\|_F^2
 &=\|\Omega_{1:K,1:K}\|_F^2+\|\Omega_{K+1,1:K}\|_F^2.
\end{align}
\end{proof}
\begin{proposition}[Polar and QR retractions]\label{prop:rigid-retract}
For $Z\in T_U\St(d,R)$,
\begin{gather}
 X:=U+Z,\quad X^\top X=I+Z^\top Z,\quad\sigma_{\min}(X)\ge1,\\
 R_U^{\mathrm{pol}}(Z):=(U+Z)(I+Z^\top Z)^{-1/2},\\
 R_U^{\mathrm{qr}}(Z):=(U+Z)\operatorname{chol}(I+Z^\top Z)^{-\top},\\
 R_U^{\mathrm{pol}}(0)=R_U^{\mathrm{qr}}(0)=U,\quad
 \mathrm DR_U^{\mathrm{pol}}(0)[Z]=\mathrm DR_U^{\mathrm{qr}}(0)[Z]=Z,\\
 \|R_U^{\mathrm{pol}}(Z)-U-Z\|_F\le
 \tfrac12\sqrt{1+\|Z\|_{\op}^2}\,\|Z\|_F^2,\\
 \|R_U^{\mathrm{qr}}(Z)-U-Z\|_F\le
 (1+\|Z\|_{\op}^2)\|Z\|_F^2.
 \label{eq:rigid-qr-error}
\end{gather}
\end{proposition}
\begin{proof}
\begin{align}
 X^\top X&=I+U^\top Z+Z^\top U+Z^\top Z=I+Z^\top Z,\\
 0&\preceq I-(I+E)^{-1/2}\preceq\tfrac12 E,\quad E:=Z^\top Z,\\
 \|R_U^{\mathrm{pol}}(Z)-X\|_F
 &\le\|X\|_{\op}\|(I+E)^{-1/2}-I\|_F
 \le\tfrac12\sqrt{1+\|Z\|_{\op}^2}\|E\|_F,\\
 L(s)L(s)^\top&=I+sE,\quad L(0)=I,\\
 L'(s)&=L(s)\mathcal T\bigl(L(s)^{-1}E L(s)^{-\top}\bigr),\\
 \mathcal T(A)&:=\operatorname{tril}(A,-1)+\tfrac12\diag(A),\\
 \|\mathcal T(A)\|_F&\le\|A\|_F,\quad
 \|L(s)^{-1}\|_{\op}\le1,\\
 \|L(1)-I\|_F&\le\int_0^1\|L'(s)\|_F\,\mathrm ds
 \le\sqrt{1+\|E\|_{\op}}\|E\|_F,\\
 \|XL(1)^{-\top}-X\|_F
 &\le\|X\|_{\op}\|L(1)^{-1}\|_{\op}\|L(1)-I\|_F\\
 &\le(1+\|E\|_{\op})\|E\|_F
 \le(1+\|Z\|_{\op}^2)\|Z\|_F^2.
\end{align}
\end{proof}
The retraction error is quadratic in the tangent step, so a small horizontal
update preserves its intended first-order motion. Comparing two different runs
additionally requires control of the QR map between their inputs. The next
lemma makes the necessary singular-value condition on that connecting segment
explicit.

\begin{lemma}[Differential of the positive-diagonal QR map]\label{lem:rigid-qr-lip}
For $X=QR$, $R_{ii}>0$, $\sigma_{\min}(X)\ge a>0$,
\begin{align}
 A&:=Q^\top(\mathrm dX)R^{-1},\\
 \mathrm dQ&=(I-QQ^\top)(\mathrm dX)R^{-1}
  +Q\bigl(\operatorname{tril}(A,-1)-\operatorname{tril}(A,-1)^\top\bigr),\\
 \|\mathrm dQ\|_F^2&\le2\|(\mathrm dX)R^{-1}\|_F^2
 \le\frac2{a^2}\|\mathrm dX\|_F^2.
\end{align}
If $\inf_{s\in[0,1]}\sigma_{\min}(X+s(Y-X))\ge a$,
\begin{equation}
 \|\operatorname{qf}(X)-\operatorname{qf}(Y)\|_F
 \le\frac{\sqrt2}{a}\|X-Y\|_F.
\end{equation}
\end{lemma}
\begin{proof}
\begin{align}
 \mathrm dX&=(\mathrm dQ)R+Q\,\mathrm dR,\\
 Q^\top(\mathrm dX)R^{-1}&=Q^\top\mathrm dQ+(\mathrm dR)R^{-1},\\
 (Q^\top\mathrm dQ)^\top&=-Q^\top\mathrm dQ,\qquad
 \operatorname{tril}((\mathrm dR)R^{-1},-1)=0,\\
 \|\operatorname{tril}(A,-1)-\operatorname{tril}(A,-1)^\top\|_F^2
 &=2\|\operatorname{tril}(A,-1)\|_F^2\le2\|A\|_F^2,\\
 \operatorname{qf}(Y)-\operatorname{qf}(X)
 &=\int_0^1\mathrm D\operatorname{qf}(X+s(Y-X))[Y-X]\,\mathrm ds.
\end{align}
\end{proof}
